% SPDX-License-Identifier: CC-BY-4.0
\documentclass[11pt,letterpaper]{article}

\usepackage[T1]{fontenc}
\usepackage{amsmath,amssymb}
\usepackage{tikz}
\usepackage{pgfplots}
\pgfplotsset{compat=1.18}
\usepackage{algorithm}
\usepackage{algpseudocode}
\usepackage{booktabs}
\usepackage{enumitem}
\usepackage[hidelinks]{hyperref}

\newcommand{\N}{\mathbb{N}}
\newcommand{\Z}{\mathbb{Z}}
\newcommand{\Q}{\mathbb{Q}}
\newcommand{\R}{\mathbb{R}}
\newcommand{\C}{\mathbb{C}}
\newcommand{\Map}[3]{#1 \colon #2 \to #3}
\newcommand{\Mod}[1]{#1\text{-}\mathrm{Mod}}
\newcommand{\GroupHom}{\mathrm{Hom}_{\mathbf{Grp}}}
\newcommand{\RingHom}{\mathrm{Hom}_{\mathbf{Ring}}}

\newtheorem{theorem}{Theorem}
\newtheorem{lemma}[theorem]{Lemma}
\newtheorem{corollary}[theorem]{Corollary}
\newtheorem{definition}[theorem]{Definition}
\newtheorem{remark}[theorem]{Remark}
\newenvironment{proof}{\noindent\textit{Proof.}}{\hfill\(\square\)}

\title{Generalized Collatz Library}
\author{}
\date{\today}

\begin{document}
\maketitle
\section*{Unit return at multiplier one}

\begin{abstract}
Let \(T_j(q)\) be the sum of \(q^0,\dots,q^{j-1}\), with the empty sum equal to
zero, and let \(A_j(q,c,n)=q^j n+c\,T_j(q)\). When the multiplier is \(1\),
the offset sum is the time itself, so \(A_j(1,c,n)=n+jc\). If \(j\) is the
first time the signed divisor \(d\) divides this state, the raw orbit agrees
with the translation through time \(j\), and the next raw step is the exact
quotient \((n+jc)/d\). Cancelling \(|d|\) leaves
\[
|d|\cdot\bigl|(n+jc)/d\bigr|\le |n|+j|c|.
\]
Since \(|1|=1<|d|\), this is a core inequality at that seed, with expansion
\(1\) and additive term \(j|c|\). The term still depends on the hitting time.
\end{abstract}

\section{The raw step and the affine state}

A parameter record \(p\) consists of integers \(q\), \(c\), and \(d\) with
\(|d|\ge 2\). In particular \(d\neq 0\). The raw step is
\[
S_p(n)=
\begin{cases}
n/d, & \text{if \(d\) divides \(n\)},\\
qn+c, & \text{otherwise.}
\end{cases}
\]
Write \(S_p^{0}(n)=n\) and \(S_p^{k+1}(n)=S_p(S_p^{k}(n))\). The offset sum is
defined by \(T_0(q)=0\) and \(T_{j+1}(q)=T_j(q)+q^j\). It satisfies
\(T_{j+1}(q)=q\,T_j(q)+1\), and the unit case \(T_j(1)=j\) is used below as a
known evaluation of this sum.

\begin{definition}
\label{def:affine-state}
The affine state at time \(j\) is \(A_j(q,c,n)=q^j n+c\,T_j(q)\).
\end{definition}

\begin{definition}
\label{def:first-hit}
Let \(p\) have multiplier \(q\), offset \(c\), and divisor \(d\). The time
\(j\) is a first hit of the seed \(n\), written \(\operatorname{Hit}(p,n,j)\),
when \(d\) divides \(A_j(q,c,n)\) and \(d\) divides \(A_i(q,c,n)\) for no
natural number \(i<j\).
\end{definition}

When \(j=0\), the failure clause is vacuous, so a first hit at time zero is
divisibility of the seed. The predicate does not assert that a first hit
exists.

\begin{theorem}
\label{thm:unit-affine}
For every \(c,n\in\Z\) and every \(j\in\N\),
\[
A_j(1,c,n)=n+jc.
\]
\end{theorem}

\begin{proof}
The unit evaluation of the offset sum gives \(T_j(1)=j\). The power \(1^j\)
equals \(1\), including the case \(j=0\), where both sides are \(1\). Therefore
\begin{align*}
A_j(1,c,n)
&=1^j\,n+c\,T_j(1)\\
&=1\cdot n+c\cdot j\\
&=n+jc,
\end{align*}
where the last step is commutativity of multiplication in \(\Z\).
\end{proof}

\section{The division-free prefix}

The next theorem is the bridge from the affine state to the raw orbit. Its
hypothesis is non-divisibility of affine states before time \(j\).
Divisibility at time \(j\) itself is not forbidden.

\begin{theorem}
\label{thm:unit-prefix}
Let \(p\) be a parameter record with multiplier \(q\), offset \(c\), and
divisor \(d\), and assume \(q=1\). Let \(n\in\Z\) and \(j\in\N\). If
\(d\nmid A_i(q,c,n)\) for every natural number \(i<j\), then
\[
S_p^{j}(n)=n+jc.
\]
\end{theorem}

\begin{proof}
The proof is by induction on \(j\).

If \(j=0\), then \(S_p^{0}(n)=n\) by definition of iteration, and
\(n+0\cdot c=n\). The hypothesis on earlier times is empty.

Suppose the statement holds at time \(j\), and fix a seed for which
\(d\nmid A_i(q,c,n)\) whenever \(i<j+1\). Every \(i<j\) satisfies \(i<j+1\),
so the inductive hypothesis applies and yields \(S_p^{j}(n)=n+jc\).
Theorem~\ref{thm:unit-affine}, together with \(q=1\), gives
\(A_j(q,c,n)=n+jc\). The hypothesis at time \(j\) says that \(d\) does not
divide \(A_j(q,c,n)\). Substituting the two equalities, \(d\) does not divide
\(S_p^{j}(n)\).

The raw step on a non-divisible state is the affine branch, so
\[
S_p^{j+1}(n)=S_p(S_p^{j}(n))=q\,S_p^{j}(n)+c.
\]
Insert \(q=1\) and the inductive evaluation:
\[
S_p^{j+1}(n)=1\cdot(n+jc)+c=n+jc+c.
\]
The natural-number successor casts to the integer successor,
\((j+1)=j+1\) in \(\Z\), hence \(n+jc+c=n+(j+1)c\), as required.
\end{proof}

\section{The return}

\begin{theorem}
\label{thm:unit-return}
Let \(p\) be a parameter record with multiplier \(q\), offset \(c\), and
divisor \(d\), and assume \(q=1\). If \(\operatorname{Hit}(p,n,j)\), then
\[
S_p\bigl(S_p^{j}(n)\bigr)=(n+jc)/d.
\]
The quotient is an integer.
\end{theorem}

\begin{proof}
Write \(\operatorname{Hit}(p,n,j)\) as the conjunction of
\(d\mid A_j(q,c,n)\) and \(d\nmid A_i(q,c,n)\) for every \(i<j\).
Theorem~\ref{thm:unit-prefix} applies to the second clause and gives
\(S_p^{j}(n)=n+jc\). Theorem~\ref{thm:unit-affine} and the hypothesis \(q=1\)
give \(A_j(q,c,n)=n+jc\). Therefore \(d\) divides \(n+jc\), and hence \(d\)
divides \(S_p^{j}(n)\).

On a divisible state the raw step is exact division by \(d\):
\[
S_p\bigl(S_p^{j}(n)\bigr)=S_p^{j}(n)/d=(n+jc)/d.
\]
Exactness is the integer identity \(d\cdot((n+jc)/d)=n+jc\), which holds
whenever \(d\) divides \(n+jc\). Admissibility \(|d|\ge 2\) is not required for
this cancellation; it is part of the record, and it is used only to place the
same quotient in the core comparison below.
\end{proof}

\section{The magnitude bound}

\begin{theorem}
\label{thm:unit-magnitude}
Let \(p\) be a parameter record with multiplier \(q\), offset \(c\), and
divisor \(d\), and assume \(q=1\). If \(\operatorname{Hit}(p,n,j)\), and if
\(r=S_p(S_p^{j}(n))\), then
\[
|d|\cdot|r|\le |n|+j|c|.
\]
\end{theorem}

\begin{proof}
Theorem~\ref{thm:unit-return} gives \(r=(n+jc)/d\). The same argument shows
that \(d\) divides \(n+jc\). For any integers \(a\) and \(b\) with \(b\)
dividing \(a\), the absolute value of the quotient in \(\Z\) is the
quotient of the absolute values in \(\N\):
\[
|a/b|=|a|/|b|.
\]
Taking \(a=n+jc\) and \(b=d\) gives \(|r|=|n+jc|/|d|\). Divisibility in
\(\Z\) implies divisibility of absolute values, so \(|d|\) divides
\(|n+jc|\) in \(\N\). Multiplying by \(|d|\) therefore cancels that
division:
\[
|d|\cdot|r|=|d|\cdot\bigl(|n+jc|/|d|\bigr)=|n+jc|.
\]
The triangle inequality in \(\Z\) gives \(|n+jc|\le |n|+|jc|\). Since \(j\) is
a natural number, \(|jc|=|j|\,|c|=j|c|\). Combining these relations yields
\[
|d|\cdot|r|=|n+jc|\le |n|+j|c|.
\]
\end{proof}

\begin{theorem}
\label{thm:unit-core}
Let \(p\) be a parameter record with multiplier \(q\), offset \(c\), and
divisor \(d\), and assume \(q=1\). If \(\operatorname{Hit}(p,n,j)\), and if
\(r=S_p(S_p^{j}(n))\), then \(|q|<|d|\) and
\[
|d|\cdot|r|\le |q|\cdot|n|+j|c|.
\]
Thus, at the seed \(n\), the return obeys the core inequality with expansion
\(A=|q|=1\) and additive term \(B=j|c|\).
\end{theorem}

\begin{proof}
The hypothesis \(q=1\) gives \(|q|=1\). The record supplies \(|d|\ge 2\), so
\(|q|<|d|\). Theorem~\ref{thm:unit-magnitude} is the inequality
\(|d|\cdot|r|\le |n|+j|c|\). Multiplying \(|n|\) by \(|q|=1\) does not change
the right-hand side, and the resulting inequality is the one stated.
\end{proof}

\begin{remark}
\label{rem:not-uniform}
The additive term \(B=j|c|\) depends on the hitting time of the seed under
consideration. The comparison \(|q|<|d|\) is uniform for every record with
multiplier \(1\), because it uses only \(|d|\ge 2\). No proof is given here
that a first hit, when it exists, satisfies \(j\le |d|-1\). Without that
clock bound, \(B\) is not replaced by \((|d|-1)|c|\) or by any other quantity
independent of \(j\).
\end{remark}

\begin{remark}
\label{rem:not-claimed}
The three comparisons do not assert that any seed has a first hit. They do not
assert the core inequality at a state that is not the seed of the chosen hit,
so they do not by themselves put every orbit into a finite invariant interval.
Nothing is claimed for a multiplier other than \(1\). Equality of the raw
iterate with \(n+jc\) uses the division-free prefix and is not claimed at any
time after the division.
\end{remark}

\end{document}
